\chapter{Finite Fields}

\begin{center}\itshape
A QR code can lose almost a third of its squares under a coffee stain and still be read.\\
How can a message repair itself?\\
The answer is polynomial arithmetic in a field with $256$ elements.
\end{center}

In 1960 Irving Reed and Gustave Solomon, two researchers at MIT's Lincoln Laboratory, published a five-page paper on ``polynomial codes over certain
finite fields''. Their codes now sit in every QR code, CD and DVD, and they carried the Voyager probes' pictures of Uranus and Neptune back across
billions of kilometres of noise. The mathematics they needed was a century old: the finite fields that \'Evariste Galois described in 1830.

The fields $\Z/p\Z$ of Volume~1 are finite. Are there others? Chapter~8 gives a way to build them: if $P$ is irreducible of degree $n$ over $\mathbb F_p=\Z/p\Z$, Kronecker's construction $\mathbb F_p[X]/(P)$ is a field with $p^n$ elements. This chapter shows that every finite field arises this way, that there is exactly one field of each prime power size, and that its multiplicative group is cyclic, as Volume~1 proved for $\Z/p\Z$. Counting the irreducible polynomials of degree $n$ gives a ``prime number theorem'' for polynomials, with an error term as good as the Riemann hypothesis would give for primes. Finally, the field with $256$ elements carries the Reed--Solomon codes of QR codes, CDs and space probes, and the AES cipher.

All rings in this chapter are commutative.

\section{The characteristic and the Frobenius map}

\begin{proposition}\label{prop:primefield}
Let $K$ be a finite field. There is a prime $p$, the \emph{characteristic} of $K$, such that $p\cdot1=0$ in $K$; the subring generated by $1$ is a field isomorphic to $\mathbb F_p=\Z/p\Z$, the \emph{prime field} of $K$; and $|K|=p^n$, where $n=[K:\mathbb F_p]$.
\end{proposition}
\begin{proof}
The map $\Z\to K$, $k\mapsto k\cdot1$, is a ring homomorphism. $K$ is finite and $\Z$ is not, so its kernel is a non-zero ideal $(p)$, $p\ge1$, and its image is isomorphic to $\Z/p\Z$ (Proposition~\ref{prop:quotring}). The image is a subring of a field, so it is a domain, and $p$ is prime (Proposition~\ref{prop:primemax}: $(p)$ is a prime ideal; $p\ne1$ because $1\ne0$ in $K$). So the image is the field $\mathbb F_p$, and $K$ is a vector space of some finite dimension $n$ over it; choosing a basis gives a bijection $K\cong\mathbb F_p^n$.
\end{proof}

\begin{lemma}[Frobenius]\label{lem:frob}
Let $K$ be a field of characteristic $p$. Then $(a+b)^p=a^p+b^p$ for all $a,b\in K$, and $\varphi:x\mapsto x^p$ is an injective ring homomorphism $K\to K$ that fixes $\mathbb F_p$. If $K$ is finite, $\varphi$ is an automorphism of $K$, the \emph{Frobenius automorphism}.
\end{lemma}
\begin{proof}
By the binomial formula, $(a+b)^p=\sum_k\binom pka^kb^{p-k}$, and for $0<k<p$ the prime $p$ divides $\binom pk=\frac{p!}{k!\,(p-k)!}$ (it divides the numerator, not the denominator). So these terms vanish in characteristic $p$. Products and $1$ are clearly respected. A ring homomorphism out of a field is injective (its kernel is an ideal not containing $1$). On $\mathbb F_p$, $x^p=x$ by Fermat's little theorem. An injective map from a finite set to itself is bijective.
\end{proof}

\begin{example}\label{ex:f4}
$X^2+X+1$ has no root in $\mathbb F_2$, so it is irreducible (degree $2$), and $\mathbb F_4=\mathbb F_2[X]/(X^2+X+1)$ is a field with $4$ elements $0,1,\alpha,\alpha+1$, where $\alpha^2=\alpha+1$. Then $\alpha^3=\alpha^2+\alpha=1$: the group $\mathbb F_4^*=\{1,\alpha,\alpha^2\}$ is cyclic of order $3$. The ring $\Z/4\Z$ also has four elements, but it is not a field ($2\cdot2=0$); a field with $p^n$ elements is \emph{not} $\Z/p^n\Z$ for $n\ge2$.
\end{example}

\section{The multiplicative group is cyclic}

\begin{theorem}\label{thm:cyclicmult}
Let $K$ be a field. Every finite subgroup $G$ of $K^*$ is cyclic. In particular, if $K$ is finite with $q$ elements, $K^*$ is cyclic of order $q-1$.
\end{theorem}
\begin{proof}
$G$ is a finite abelian group, of order $m$ say, so $G\cong\Z/d_1\oplus\dots\oplus\Z/d_r$ with $d_1\mid\dots\mid d_r$ (Corollary~\ref{cor:abelian}). Every element $x$ of $\bigoplus\Z/d_i$ satisfies $d_rx=0$, so every $x\in G$ satisfies $x^{d_r}=1$. The polynomial $X^{d_r}-1$ has at most $d_r$ roots in the field $K$, so $m\le d_r$. As $m=d_1\cdots d_r$, $r=1$ and $G\cong\Z/m$.
\end{proof}

A generator of $\mathbb F_q^*$ is a \emph{primitive element}. If $\gamma$ is one, every non-zero element is a power of $\gamma$, so $\mathbb F_q=\mathbb F_p(\gamma)\cong\mathbb F_p[X]/(m_\gamma)$ (Theorem~\ref{thm:simple}): every finite field is obtained by Kronecker's construction. Multiplication in $\mathbb F_q^*$ is addition of exponents modulo $q-1$, as on the circle of Figure~\ref{fig:f16}; addition is addition of coordinates in $\mathbb F_p^n$. A finite field is the meeting point of these two very different pictures.

\begin{figure}[ht]
\centering
\fig{ch09-f16}
\caption{Left: $\mathbb F_{16}=\mathbb F_2[\alpha]$ with $\alpha^4=\alpha+1$. Its $15$ non-zero elements are the powers of $\alpha$, placed on a circle, so multiplying turns the circle; under each power are its coordinates on $\alpha^3,\alpha^2,\alpha,1$, which are added bit by bit (exclusive or). Right: the number $N_2(n)$ of irreducible polynomials of degree $n$ over $\mathbb F_2$ follows $2^n/n$ closely (Theorem~\ref{thm:countirr}); the vertical scale is logarithmic.}
\label{fig:f16}
\end{figure}

\section{Existence and uniqueness}

\begin{lemma}\label{lem:xqx}
If $K$ is a field with $q$ elements, then $x^q=x$ for every $x\in K$, and
\[
X^q-X=\prod_{a\in K}(X-a)\quad\text{in }K[X].
\]
\end{lemma}
\begin{proof}
For $x\ne0$, $x^{q-1}=1$ by Lagrange's theorem in the group $K^*$ of order $q-1$; multiply by $x$. So the $q$ elements of $K$ are roots of the monic polynomial $X^q-X$ of degree $q$, which is therefore their product $\prod(X-a)$ (each root splits off a factor, Volume~0).
\end{proof}

\begin{theorem}[Galois, Moore]\label{thm:fq}
For every prime $p$ and $n\ge1$ there is a field with $q=p^n$ elements, and any two such fields are isomorphic: it is the splitting field of $X^q-X$ over $\mathbb F_p$. We denote it by $\mathbb F_q$.
\end{theorem}
\begin{proof}
\emph{Existence.} Let $L$ be a splitting field of $P=X^q-X$ over $\mathbb F_p$ (Proposition~\ref{prop:splitexists}). The derivative is $P'=qX^{q-1}-1=-1$ (as $q=0$ in $L$), so $P$ has no repeated root: at a double root $\alpha$ we would have $P'(\alpha)=0$, as in Lemma~\ref{lem:sep}. So $P$ has exactly $q$ roots in $L$. They are the fixed points of $\varphi^n:x\mapsto x^{p^n}$, a ring homomorphism (Lemma~\ref{lem:frob}), so they form a subfield of $L$: if $x^q=x$ and $y^q=y$, then $(x\pm y)^q=x\pm y$, $(xy)^q=xy$ and $(x^{-1})^q=x^{-1}$. This subfield contains $\mathbb F_p$ and all the roots of $P$, which generate $L$, so it is $L$, and $|L|=q$.

\emph{Uniqueness.} Let $K$ have $q$ elements. Its prime field is isomorphic to $\mathbb F_p$, and by Lemma~\ref{lem:xqx}, $X^q-X$ splits in $K$ with the elements of $K$ as roots, which certainly generate $K$. So $K$ is a splitting field of $X^q-X$ over $\mathbb F_p$, and two splitting fields are isomorphic (Theorem~\ref{thm:splitunique}).
\end{proof}

\section{Subfields and the Galois group}

We use twice that if $d\mid n$, then $a^d-1$ divides $a^n-1$, for integers and for polynomials alike: with $n=de$, $a^n-1=(a^d-1)(a^{d(e-1)}+\dots+a^d+1)$.

\begin{theorem}\label{thm:subfields}
The subfields of $\mathbb F_{p^n}$ are the sets $\{x:x^{p^d}=x\}$ for $d\mid n$; there is exactly one subfield with $p^d$ elements for each divisor $d$ of $n$, and no other.
\end{theorem}
\begin{proof}
Let $K$ be a subfield. It contains the prime field, so $|K|=p^d$ with $d=[K:\mathbb F_p]$, and the tower law gives $n=d\,[\mathbb F_{p^n}:K]$, so $d\mid n$. By Lemma~\ref{lem:xqx}, $K$ is contained in the set of roots of $X^{p^d}-X$, which has at most $p^d$ elements, so $K$ equals that set.

Conversely let $d\mid n$. Then $p^d-1\mid p^n-1$, hence $X^{p^d-1}-1$ divides $X^{p^n-1}-1$, and $X^{p^d}-X$ divides $X^{p^n}-X$. The latter splits in $\mathbb F_{p^n}$ with distinct roots (Lemma~\ref{lem:xqx}), so the former has $p^d$ distinct roots there. They are the fixed points of the homomorphism $\varphi^d$, so they form a subfield, with $p^d$ elements.
\end{proof}

So the subfields of $\mathbb F_{p^n}$ form a copy of the lattice of divisors of $n$: $\mathbb F_{2^{12}}$ contains $\mathbb F_{2^4}$ and $\mathbb F_{2^6}$, which meet in $\mathbb F_{2^2}$, but $\mathbb F_8\not\subseteq\mathbb F_{16}$.

\begin{theorem}\label{thm:galfq}
The automorphisms of $\mathbb F_{p^n}$ form a cyclic group of order $n$, generated by the Frobenius automorphism $\varphi$.
\end{theorem}
\begin{proof}
For $1\le k<n$, $\varphi^k\ne\mathrm{id}$: the polynomial $X^{p^k}-X$ has at most $p^k<p^n$ roots, so some $x$ has $x^{p^k}\ne x$. And $\varphi^n=\mathrm{id}$ by Lemma~\ref{lem:xqx}. So $\langle\varphi\rangle$ has order $n$. Every automorphism fixes $1$, hence the prime field. $\mathbb F_{p^n}$ is a splitting field of $X^{p^n}-X$ over $\mathbb F_p$, a polynomial with distinct roots, so by Theorem~\ref{thm:splitunique} it has exactly $[\mathbb F_{p^n}:\mathbb F_p]=n$ automorphisms. They are the powers of $\varphi$.
\end{proof}

This is the first complete example of the Galois correspondence of Chapter~10: the subgroups of $\Z/n\Z$ (one for each $d\mid n$) match the subfields (one for each $d\mid n$), the subgroup $\langle\varphi^d\rangle$ fixing exactly $\mathbb F_{p^d}$.

\section{Counting irreducible polynomials}

\begin{think}
How many monic irreducible polynomials of degree $2$ are there over $\mathbb F_2$? Over $\mathbb F_3$? (A monic quadratic is reducible iff it has a root.) Can you guess a formula over $\mathbb F_q$?
\end{think}

Everything above works over any finite field $\mathbb F_q$ in place of $\mathbb F_p$, as $\mathbb F_q\subseteq\mathbb F_{q^n}$ by Theorem~\ref{thm:subfields} (if $q=p^m$, then $m\mid mn$), and $x\mapsto x^q=\varphi^m(x)$ is a homomorphism fixing exactly $\mathbb F_q$.

\begin{theorem}\label{thm:xqnx}
In $\mathbb F_q[X]$, $X^{q^n}-X$ is the product of all monic irreducible polynomials whose degree divides $n$.
\end{theorem}
\begin{proof}
The derivative of $X^{q^n}-X$ is $-1$, so it has no repeated factor: if $P^2$ divided it, every root of $P$ in a splitting field would be a double root. It is monic, and $\mathbb F_q[X]$ is a UFD, so it suffices to show that a monic irreducible $P$, of degree $d$, divides $X^{q^n}-X$ iff $d\mid n$. Let $K=\mathbb F_q[X]/(P)$, a field with $q^d$ elements, and $\alpha$ the class of $X$, whose minimal polynomial is $P$. For $k\ge0$ let $\psi_k(x)=x^{q^k}$, a homomorphism of $K$, and $S=\{k\ge0:\psi_k=\mathrm{id}_K\}$. By Lemma~\ref{lem:xqx}, $d\in S$.

If $d\mid n$, then $\psi_n=(\psi_d)^{n/d}=\mathrm{id}$, so $\alpha^{q^n}=\alpha$, and $P=m_\alpha$ divides $X^{q^n}-X$.

Conversely, suppose $P\mid X^{q^n}-X$, so $\psi_n(\alpha)=\alpha$. The fixed points of $\psi_n$ form a subfield containing $\mathbb F_q$ and $\alpha$, which generate $K$; so $\psi_n=\mathrm{id}$ and $n\in S$. If $k\ge l$ are in $S$, then $\psi_k=\psi_{k-l}\circ\psi_l=\psi_{k-l}$, so $k-l\in S$. Running Euclid's algorithm by subtractions, $g=\gcd(n,d)\in S$: every one of the $q^d$ elements of $K$ is a root of $X^{q^g}-X$, so $q^d\le q^g$, $d\le g$, and $d=g$ divides $n$.
\end{proof}

\begin{theorem}[Gauss]\label{thm:countirr}
Let $N_q(n)$ be the number of monic irreducible polynomials of degree $n$ over $\mathbb F_q$. Then
\[
\sum_{d\mid n}d\,N_q(d)=q^n,\qquad N_q(n)=\frac1n\sum_{d\mid n}\mu(d)\,q^{n/d},
\]
and
\[
\Bigl|N_q(n)-\frac{q^n}n\Bigr|<\frac{2q^{n/2}}n.
\]
In particular $N_q(n)\ge1$ for every $n\ge1$.
\end{theorem}
\begin{proof}
Compare degrees in Theorem~\ref{thm:xqnx}: the irreducibles of degree $d$, for each $d\mid n$, contribute $d\,N_q(d)$. Möbius inversion (Volume~1, Chapter~4), applied to $f(d)=dN_q(d)$ and $g(n)=q^n$, gives the formula. In it, the term $d=1$ is $q^n$, and the others have $n/d\le n/2$, so with $m=\lfloor n/2\rfloor$,
\[
\begin{aligned}\Bigl|nN_q(n)-q^n\Bigr|&\le\sum_{k=1}^mq^k=\frac{q(q^m-1)}{q-1}\eqbr<\frac q{q-1}q^m\le2q^{n/2}.\end{aligned}
\]
For the last claim: $N_q(1)=q$, and for $n\ge2$, $2q^{n/2}\le q^{n/2+1}\le q^n$, so $nN_q(n)>0$.
\end{proof}

So there are $(q^2-q)/2$ irreducible monic quadratics over $\mathbb F_q$: one over $\mathbb F_2$ ($X^2+X+1$) and three over $\mathbb F_3$ ($X^2+1$, $X^2+X+2$, $X^2+2X+2$). Directly: there are $q^2$ monic quadratics, of which $q$ are squares $(X-a)^2$ and $\binom q2$ are products $(X-a)(X-b)$ with $a\ne b$; $q^2-q-\frac{q(q-1)}2=\frac{q^2-q}2$.

\begin{remark}
Think of $X^{q^n}$, the number of monic polynomials of degree $n$, as ``$x$''. Then $n=\log_qx$, and Theorem~\ref{thm:countirr} says that the number of irreducibles of degree $n$ is $\frac x{\log_qx}$ up to an error of size $\sqrt x$: a prime number theorem for $\mathbb F_q[X]$, with an error term as small as the one the Riemann hypothesis would give for the primes (up to a logarithm). For polynomials the proof took a page; for primes it is the most famous open problem in mathematics, met again later in the series. Weil's theorem on curves over finite fields is a deep extension of this ``Riemann hypothesis over finite fields''.
\end{remark}

\section{Correcting errors: Reed--Solomon codes}

To send $k$ symbols $m_0,\dots,m_{k-1}\in\mathbb F_q$ over a noisy channel, encode them as the polynomial $f=\sum m_jX^j$ of degree $<k$, choose $n$ distinct points $a_1,\dots,a_n\in\mathbb F_q$ ($k\le n\le q$), and send the \emph{codeword} $(f(a_1),\dots,f(a_n))$. The \emph{Hamming distance} between two words of $\mathbb F_q^n$ is the number of positions where they differ; it satisfies the triangle inequality, since if $u_i\ne w_i$ then $u_i\ne v_i$ or $v_i\ne w_i$.

\begin{proposition}\label{prop:rsdist}
Two different codewords are at distance at least $n-k+1$. Consequently, if at most $t=\lfloor\frac{n-k}2\rfloor$ symbols are corrupted, the sent codeword is the unique codeword closest to the received word.
\end{proposition}
\begin{proof}
If $f\ne g$ have degree $<k$, then $f-g$ is a non-zero polynomial of degree $<k$, so it has at most $k-1$ roots: $f(a_i)=g(a_i)$ for at most $k-1$ indices $i$. If the sent word $c$ and the received word $y$ are at distance $e\le t$, and $c'\ne c$ is another codeword, then $d(c',y)\ge d(c,c')-d(c,y)\ge n-k+1-t>t\ge e$, because $2t\le n-k$.
\end{proof}

Finding the closest codeword by trying all $q^k$ messages is hopeless when $q=256$ and $k=200$. Berlekamp and Welch found a way with linear algebra: an unknown \emph{error locator} polynomial $E$ vanishes at the bad positions, and $fE$ is then determined by a linear system.

\begin{theorem}[Berlekamp--Welch decoding]\label{thm:bw}
Let $y\in\mathbb F_q^n$ differ from the codeword of $f$ in at most $t=\lfloor\frac{n-k}2\rfloor$ positions. Then the linear system
\[
Q(a_i)=y_i\,E(a_i)\qquad(1\le i\le n),
\]
in the unknowns $E$ monic of degree $t$ and $Q$ of degree $<k+t$, has a solution, and every solution satisfies $Q=fE$. So $f$ is the quotient of $Q$ by $E$.
\end{theorem}
\begin{proof}
\emph{Existence.} Let $B$ be the set of the $e\le t$ positions where $y_i\ne f(a_i)$, and $E=X^{t-e}\prod_{i\in B}(X-a_i)$, $Q=fE$. Then $E$ is monic of degree $t$, $\deg Q<k+t$; for $i\in B$ both sides vanish, and for $i\notin B$, $Q(a_i)=f(a_i)E(a_i)=y_iE(a_i)$.

\emph{Every solution works.} Let $(E,Q)$ be any solution and $R=Q-fE$, of degree $<k+t$. For each of the at least $n-t$ indices $i\notin B$, $R(a_i)=y_iE(a_i)-f(a_i)E(a_i)=0$. As $n-t\ge k+t$ (because $2t\le n-k$), $R$ has more roots than its degree, so $R=0$.
\end{proof}

The system has $n$ equations and $k+2t\le n$ unknowns (the $t$ lower coefficients of $E$ and the $k+t$ coefficients of $Q$); Gaussian elimination solves it in $O(n^3)$ operations in $\mathbb F_q$, and faster algorithms (Berlekamp--Massey, Euclid's algorithm on polynomials) need $O(n^2)$. QR codes, CDs, DVDs, RAID~6 disk arrays and the Voyager probes all use Reed--Solomon codes over $\mathbb F_{256}$, whose elements are bytes; the highest QR level adds enough redundancy to recover about $30\%$ of the data. The same field, built as $\mathbb F_2[X]/(X^8+X^4+X^3+X+1)$, is the arithmetic of the AES cipher.

\begin{beauty}
A finite field is completely described by two numbers, $p$ and $n$, and yet it carries a rich structure: a cyclic multiplicative group, a cyclic Galois group generated by $x\mapsto x^p$, a lattice of subfields copying the divisors of $n$, and a count of irreducible polynomials that mirrors the distribution of primes. Galois discovered these ``imaginary'' fields in 1830, the year of his work on equations. Today they are everywhere in computing: error-correcting codes, AES, elliptic-curve cryptography, hashing, pseudo-random generators. The Frobenius automorphism, in turn, becomes in Volume~12 the bridge between primes and Galois groups (Frobenius elements, and Chebotarev's theorem).
\end{beauty}

\section{Exercises}

\begin{exercise}[$\star$]\label{ex:f4tables}
Write the addition and multiplication tables of $\mathbb F_4$ (Example~\ref{ex:f4}). Show that $\mathbb F_4$ and $\Z/4\Z$ are not even isomorphic as additive groups, and that $\mathbb F_2[X]/(X^2+1)$ is not a field.
\end{exercise}

\begin{exercise}[$\star$]\label{ex:irrF2}
List the monic irreducible polynomials of degree $\le4$ over $\mathbb F_2$, check the count of Theorem~\ref{thm:countirr}, and factor $X^{16}-X$ in $\mathbb F_2[X]$.
\end{exercise}

\begin{exercise}[$\star\star$]\label{ex:primpoly}
An irreducible $P\in\mathbb F_p[X]$ is \emph{primitive} if the class of $X$ generates $(\mathbb F_p[X]/(P))^*$. Show that $X^4+X+1$ is primitive over $\mathbb F_2$, but $X^4+X^3+X^2+X+1$ is not. How many primitive elements does $\mathbb F_{16}$ have, and how many primitive polynomials of degree $4$ are there over $\mathbb F_2$?
\end{exercise}

\begin{exercise}[$\star\star$]\label{ex:minus1}
Let $q$ be odd. Show that $-1$ is a square in $\mathbb F_q$ iff $q\equiv1\pmod4$. Deduce that $-1$ is a square in $\mathbb F_9$, and find a square root of $-1$ in $\mathbb F_9=\mathbb F_3[X]/(X^2+X+2)$.
\end{exercise}

\begin{exercise}[$\star\star$]\label{ex:twosq}
Show that every element of a finite field is a sum of two squares.
\end{exercise}

\begin{exercise}[$\star\star$]\label{ex:powersums}
\begin{enumerate}[label=(\alph*)]
\item Show that for an integer $k\ge0$, $\displaystyle\sum_{x\in\mathbb F_q}x^k$ is $-1$ if $k\ge1$ and $(q-1)\mid k$, and $0$ otherwise (with $0^0=1$).
\item (Wilson) Show that the product of all non-zero elements of $\mathbb F_q$ is $-1$.
\end{enumerate}
\end{exercise}

\begin{exercise}[$\star\star$]\label{ex:rs7}
Over $\mathbb F_7$, a message $f$ of degree $<2$ is sent as $(f(1),f(2),\dots,f(6))$, and $y=(5,3,2,4,0,1)$ is received. Assuming at most two errors, decode it with Theorem~\ref{thm:bw}.
\end{exercise}

\begin{exercise}[$\star\star\star$]\label{ex:chevalley}
(Chevalley--Warning.) Let $F\in\mathbb F_q[X_1,\dots,X_n]$ have total degree $d<n$. Show that the number $N$ of zeros of $F$ in $\mathbb F_q^n$ is divisible by $p$, the characteristic. Deduce that if $F(0)=0$, then $F$ has a non-zero zero; for instance, every quadratic form in $3$ variables over $\mathbb F_q$ has a non-trivial zero.
\end{exercise}

\section{Hints}

\begin{hint}{ex:f4tables}
In $\mathbb F_4$, $x+x=0$ for every $x$. In $\mathbb F_2[X]$, $X^2+1=(X+1)^2$.
\end{hint}
\begin{hint}{ex:irrF2}
A polynomial of degree $2$ or $3$ is irreducible iff it has no root; one of degree $4$ iff it has no root and is not $(X^2+X+1)^2$.
\end{hint}
\begin{hint}{ex:primpoly}
The order of $X$ divides $15$; compute $X^3$ and $X^5$. For the counts, each primitive polynomial of degree $4$ has $4$ roots, all primitive.
\end{hint}
\begin{hint}{ex:minus1}
$-1$ is the unique element of order $2$ in the cyclic group $\mathbb F_q^*$. Try elements $a+b\alpha$ with $\alpha^2=-\alpha-2$.
\end{hint}
\begin{hint}{ex:twosq}
For $q$ odd, how many squares are there? Count the values of $x^2$ and of $a-y^2$.
\end{hint}
\begin{hint}{ex:powersums}
(a) Write the non-zero elements as powers of a primitive element: a geometric sum. (b) $X^{q-1}-1=\prod_{a\ne0}(X-a)$; look at the constant term.
\end{hint}
\begin{hint}{ex:rs7}
Here $n=6$, $k=2$, $t=2$: unknowns $E=X^2+e_1X+e_0$ and $Q$ of degree $\le3$. If you guess where the errors are, the existence part of the proof tells you $E$ and $Q$.
\end{hint}
\begin{hint}{ex:chevalley}
In $\mathbb F_q$, $1-F(x)^{q-1}$ is $1$ if $F(x)=0$ and $0$ otherwise. Sum over $x\in\mathbb F_q^n$, expand $F^{q-1}$ into monomials and use Exercise~\ref{ex:powersums}(a).
\end{hint}

\section{Solutions}

\begin{solution}{ex:f4tables}
\idea Compute with $\alpha^2=\alpha+1$ and $2=0$; compare the orders of elements in the two additive groups.
\writeup Writing $\beta=\alpha+1=\alpha^2$: addition is $x+x=0$, $1+\alpha=\beta$, $1+\beta=\alpha$, $\alpha+\beta=1$ (and $0$ is neutral); multiplication is $\alpha\cdot\alpha=\beta$, $\alpha\beta=\alpha^3=1$, $\beta\beta=\alpha^4=\alpha$ (and $1$ is neutral, $0$ absorbing). In $(\mathbb F_4,+)$ every element satisfies $2x=0$, so it is $(\Z/2\Z)^2$; in $\Z/4\Z$, $1$ has order $4$. So they are not isomorphic groups, and in fact the additive group of any field of characteristic $p$ has exponent $p$. Finally $X^2+1=(X+1)^2$ in $\mathbb F_2[X]$, so the class of $X+1$ is a non-zero element of $\mathbb F_2[X]/(X^2+1)$ whose square is $0$: the ring is not a domain.
\end{solution}

\begin{solution}{ex:irrF2}
\idea Low degree: irreducible means no factor of degree $\le$ half the degree; then check against Theorem~\ref{thm:xqnx}.
\writeup A polynomial has the root $0$ iff its constant term is $0$, and the root $1$ iff it has an even number of terms. Degree $1$: $X$, $X+1$. Degree $2$: $X^2+X+1$ only. Degree $3$ (no root): $X^3+X+1$, $X^3+X^2+1$. Degree $4$: the candidates with no root are $X^4+X+1$, $X^4+X^3+1$, $X^4+X^3+X^2+X+1$ and $X^4+X^2+1=(X^2+X+1)^2$, so three irreducibles. The counts $2,1,2,3$ agree with $N_2(n)=\frac1n\sum_{d\mid n}\mu(d)2^{n/d}$: $2$, $\frac{4-2}2=1$, $\frac{8-2}3=2$, $\frac{16-4}4=3$. By Theorem~\ref{thm:xqnx}, $X^{16}-X$ is the product of the irreducibles of degree $1$, $2$, $4$:
\[
\begin{aligned}
X^{16}-X&=X(X+1)(X^2+X+1)\eqbr{}\cdot(X^4+X+1)(X^4+X^3+1)\eqbr{}\cdot(X^4+X^3+X^2+X+1),
\end{aligned}
\]
of degree $1+1+2+12=16$.
\end{solution}

\begin{solution}{ex:primpoly}
\idea In a cyclic group of order $15$, an element is a generator iff its order is not $1$, $3$ or $5$.
\writeup For $P=X^4+X+1$: it is irreducible (Exercise~\ref{ex:irrF2}), so $K=\mathbb F_2[X]/(P)$ is $\mathbb F_{16}$ and the order of $x$ divides $15$. Now $x^3\ne1$ (the classes $1,x,x^2,x^3$ are a basis) and $x^5=x\cdot x^4=x(x+1)=x^2+x\ne1$. So $x$ has order $15$ (this is Figure~\ref{fig:f16}). For $P=X^4+X^3+X^2+X+1$: $(X-1)P=X^5-1$, so $x^5=1$ and $x$ has order $5$; $P$ is irreducible but not primitive.

The cyclic group $\mathbb F_{16}^*$ of order $15$ has $\varphi(15)=8$ generators. A primitive polynomial of degree $4$ is the minimal polynomial of a primitive element; conversely the minimal polynomial of a primitive element $\gamma$ has degree $4$ ($\mathbb F_2(\gamma)=\mathbb F_{16}$), and its roots are $\gamma,\gamma^2,\gamma^4,\gamma^8$ (the images by Frobenius, all distinct and all primitive, since squaring is an automorphism of $\mathbb F_{16}^*$). So the $8$ primitive elements fall into $8/4=2$ groups: there are two primitive polynomials, $X^4+X+1$ and $X^4+X^3+1$.
\end{solution}

\begin{solution}{ex:minus1}
\idea In a cyclic group, a given element is a square iff its ``exponent'' with respect to a generator is even.
\writeup Let $\gamma$ generate $\mathbb F_q^*$, of even order $q-1$. The element $\gamma^{(q-1)/2}$ has order $2$, and it is the only one, as $X^2-1$ has the two roots $\pm1$; so $-1=\gamma^{(q-1)/2}$. The squares in $\mathbb F_q^*$ are the $(\gamma^j)^2=\gamma^{2j}$; as $q-1$ is even, the parity of an exponent $i$ of $\gamma^i$ is well defined modulo $q-1$, so $\gamma^i$ is a square iff $i$ is even. So $-1$ is a square iff $(q-1)/2$ is even, iff $q\equiv1\pmod4$. As $9\equiv1\pmod4$, $-1$ is a square in $\mathbb F_9$. With $\alpha^2=-\alpha-2=2\alpha+1$: $(\alpha+2)^2=\alpha^2+4\alpha+4=\alpha^2+\alpha+1=3\alpha+2=2=-1$. So $\alpha+2$ is a square root of $-1$ (and so is $-(\alpha+2)=2\alpha+1$).
\end{solution}

\begin{solution}{ex:twosq}
\idea Pigeonhole: two sets with more than $q/2$ elements each must meet.
\writeup If $q$ is even, Frobenius $x\mapsto x^2$ is a bijection, so every $a$ is a square, $a=x^2+0^2$. If $q$ is odd, $x\mapsto x^2$ on $\mathbb F_q^*$ is two-to-one ($x^2=y^2$ iff $y=\pm x$, and $x\ne-x$), so there are $\frac{q-1}2+1=\frac{q+1}2$ squares, including $0$. For $a\in\mathbb F_q$, the sets $\{x^2\}$ and $\{a-y^2\}$ both have $\frac{q+1}2$ elements; their sizes add up to $q+1>q$, so they meet: $x^2=a-y^2$.
\end{solution}

\begin{solution}{ex:powersums}
\idea Sum over a cyclic group of a character-like power: either everything is $1$, or the sum is a geometric series that vanishes.
\writeup (a) If $k=0$ the sum is $q\cdot1=0$. Let $k\ge1$; then $0^k=0$, and with $\gamma$ a primitive element, $\sum_{x}x^k=\sum_{j=0}^{q-2}(\gamma^k)^j$. If $(q-1)\mid k$, $\gamma^k=1$ and the sum is $q-1=-1$. Otherwise $\gamma^k\ne1$ and the geometric sum is $\frac{(\gamma^k)^{q-1}-1}{\gamma^k-1}=0$.

(b) By Lemma~\ref{lem:xqx} divided by $X$, $X^{q-1}-1=\prod_{a\ne0}(X-a)$. Comparing constant terms, $-1=(-1)^{q-1}\prod_{a\ne0}a$. If $q$ is odd, $(-1)^{q-1}=1$; if $q$ is even, $-1=1$. Either way $\prod_{a\ne0}a=-1$. (For $q=p$ this is Wilson's theorem $(p-1)!\equiv-1\pmod p$.)
\end{solution}

\begin{solution}{ex:rs7}
\idea Write the six linear equations $Q(a_i)=y_iE(a_i)$ and solve; or guess the error positions and check, since Theorem~\ref{thm:bw} says that any solution gives the right answer.
\writeup The unknowns are $e_0,e_1$ and $Q=q_0+q_1X+q_2X^2+q_3X^3$. Solving the $6\times6$ system modulo $7$ by elimination (the matrix is invertible) gives $E=X^2+3$ and $Q=2X^3+3X^2+6X+2$. Check, with $f=Q/E$: indeed $(2X+3)(X^2+3)=2X^3+3X^2+6X+9$, and $9=2$ in $\mathbb F_7$. The table of values at $a=1,\dots,6$ confirms the six equations:
\begin{center}
\small
\begin{tabular}{c|cccccc}
$a$ & $1$ & $2$ & $3$ & $4$ & $5$ & $6$\\\hline
$y$ & $5$ & $3$ & $2$ & $4$ & $0$ & $1$\\
$E(a)$ & $4$ & $0$ & $5$ & $5$ & $0$ & $4$\\
$yE(a)=Q(a)$ & $6$ & $0$ & $3$ & $6$ & $0$ & $4$\\
$f(a)=2a+3$ & $5$ & $0$ & $2$ & $4$ & $6$ & $1$
\end{tabular}
\end{center}
So the message is $f=3+2X$. The error locator $E=X^2+3=(X-2)(X-5)$ shows where the errors were: positions $2$ and $5$, where $3$ and $0$ were received instead of $0$ and $6$.
\end{solution}

\begin{solution}{ex:chevalley}
\idea Count zeros inside $\mathbb F_q$ with an indicator polynomial, and show that the polynomial is too small in degree to have a non-zero sum.
\writeup By Lemma~\ref{lem:xqx}, $F(x)^{q-1}=1$ if $F(x)\ne0$, so $1-F(x)^{q-1}$ is the indicator of $F(x)=0$, and in $\mathbb F_q$,
\[
\begin{aligned}N\cdot1&=\sum_{x\in\mathbb F_q^n}\bigl(1-F(x)^{q-1}\bigr)\eqbr=q^n\cdot1-\sum_{x}F(x)^{q-1}=-\sum_xF(x)^{q-1}.\end{aligned}
\]
Expand $F^{q-1}$ into monomials $cX_1^{k_1}\cdots X_n^{k_n}$ of total degree $k_1+\dots+k_n\le d(q-1)<n(q-1)$. For each, some $k_i<q-1$, and
\[
\sum_{x\in\mathbb F_q^n}x_1^{k_1}\cdots x_n^{k_n}=\prod_{j=1}^n\Bigl(\sum_{x_j\in\mathbb F_q}x_j^{k_j}\Bigr)
\]
has the factor $\sum_{x_i}x_i^{k_i}=0$ (Exercise~\ref{ex:powersums}(a): $k_i=0$, or $k_i\ge1$ not divisible by $q-1$). So $N\cdot1=0$ in $\mathbb F_q$, i.e.\ $p\mid N$. If $F(0)=0$, then $N\ge1$, so $N\ge p\ge2$, and there is a zero other than $0$. A quadratic form in $3$ variables has degree $2<3$ and vanishes at $0$.
\end{solution}

\section*{Sources and further reading}

\begin{itemize}
\item \'E.~Galois, \emph{Sur la théorie des nombres}, Bulletin des sciences mathématiques de Férussac 13 (1830), 428--435. The first construction of the fields $\mathbb F_{p^n}$, with ``imaginary roots'' of irreducible congruences.
\item E.~H.~Moore, \emph{A doubly-infinite system of simple groups}, Bull.\ New York Math.\ Soc.\ 3 (1893), 73--78. Every finite field is one of Galois' fields.
\item C.~F.~Gauss, \emph{Disquisitiones generales de congruentiis} (the unpublished Section~VIII of the \emph{Disquisitiones}, printed in his Werke, vol.~2, 1863). The count of irreducible polynomials modulo $p$.
\item I.~S.~Reed and G.~Solomon, \emph{Polynomial codes over certain finite fields}, J.\ SIAM 8 (1960), 300--304. The codes of this chapter, in five pages.
\item L.~R.~Welch and E.~R.~Berlekamp, \emph{Error correction for algebraic block codes}, US Patent 4\,633\,470 (1986). The decoding algorithm of Theorem~\ref{thm:bw}.
\item C.~Chevalley, \emph{Démonstration d'une hypothèse de M.~Artin}, Abh.\ Math.\ Sem.\ Hamburg 11 (1935), 73--75; E.~Warning, \emph{Bemerkung zur vorstehenden Arbeit von Herrn Chevalley}, same volume, 76--83. Exercise~\ref{ex:chevalley}.
\item R.~Lidl and H.~Niederreiter, \emph{Introduction to Finite Fields and Their Applications}, revised ed., Cambridge University Press, 1994. Everything about finite fields, with codes and cryptography.
\item J.~Daemen and V.~Rijmen, \emph{The Design of Rijndael: AES, the Advanced Encryption Standard}, Springer, 2002. Why AES computes in $\mathbb F_{256}$, an application of finite-field arithmetic.
\item I.~S.~Reed and G.~Solomon, \emph{Polynomial codes over certain finite fields}, J.\ SIAM 8 (1960), 300--304; S.~B.~Wicker and V.~K.~Bhargava (eds.), \emph{Reed--Solomon Codes and Their Applications}, IEEE Press, 1994, on their use by the Voyager probes and in compact discs.
\end{itemize}
